\Def Линейным рекуррентным соотношением с постоянными коэффициентами (ЛРСПК) порядка n назовем соотношение вида $a_k = \sum\limits_{i=1}^n c_{i} a_{k-i}\,+\,f(k,a_{k-1},\dots,a_{k-n})$ с заданными $a_1, \dots , a_n$.

ЛРСПК называется однородным, если $f\, \equiv \,0$

\subsection*{Нахождение k-го члена однородной ЛРСПК n-го порядка за $O(n^3\,log\,k)$}

Запишем наши первые n членов последовательности в столбик $A_0 = \begin{pmatrix}
a_{n - 1}\\
a_{n - 2}\\
\vdots \\
a_0
\end{pmatrix}$

Так же выпишем следующую матрицу перехода:
$T = \begin{pmatrix}
c_1 & c_2 & c_3 & \cdots & c_{n - 1} & c_n\\
1 & 0 & 0 & \cdots & 0 & 0\\
0 & 1 & 0 & \cdots & 0 & 0\\
\vdots & \vdots & \vdots & ~ & \vdots & \vdots \\
0 & 0 & 0 & \cdots & 1 & 0\\
\end{pmatrix}$

Заметим, что умножив $T$ на $A_0$, мы получим столбик $A_1$ следующего вида: 
$A_1 = T \cdot A_0 = \begin{pmatrix}
a_{n}\\
a_{n - 1}\\
\vdots \\
a_1
\end{pmatrix}$
Аналогично, домножив $T$ на $A_1$, получим 
$A_2 = T \cdot A_1 = \begin{pmatrix}
a_{n + 1}\\
a_{n}\\
\vdots \\
a_2
\end{pmatrix}$

Продолжая так, для любого целого неотрицательного $i$, мы получим столбик $A_i$, состоящий из $n$ подряд идущих членов последовательности, начиная с $a_i$. Пользуясь ассоциативностью произведения матриц, можно записать, что $A_i = T^i \cdot A_0$. Из этого соотношения вытекает алгоритм вычисления произвольного $a_k$:

\begin{enumerate}
    \item Инициализировать матрицы $A_0$ и $T$
    \item Возвести матрицу $T$ в степень $k$
    \item Посчитать $A_k$ как $T^k \cdot A_0$ и взять из него $a_k$
\end{enumerate}

Используя быстрое возведение в степень во втором пункте, мы будем тратить $O(n^3 \cdot \log k)$ времени. Умножение же в третьем пункте выполняется за $O(n^3)$.

Итого мы получили алгоритм, работающий за $O(n^3 \cdot \log k)$.

\pagebreak